In this section we fully tackle the problem of realisability of sets
of scenarios as expressions, that is we account also for
constraints. We begin with an example.
%The members of ${\mathcal{S}}$, in all the examples so far,
%include only default constraints.
%It turns out that in the presence of non-default constraints, when
%${\mathcal{S}}$ is not realisable by a single DTS, reasoning
%based on K-rankings alone might not provide us with enough
%information to properly partition {\mathcal{S}}$.
%----------------
% try3-time
\begin{example}
\label{ex-with-time-1}
% Consider the set of scenarios
Let $\mathcal{S}=\{\SeqS_1,\SeqS_2,\SeqS_3\}$ over $\Sigma=\{a,b,c,d\}$, where
$\SeqS_1=(cdab,\{\TT{c}{b} \leq 3\})$,
$\SeqS_2=(cadb,\{\TT{c}{d} \leq 1,\TT{c}{b} \leq 3\})$ and
$\SeqS_3=(cabd,\{\TT{c}{d} \leq 1\})$.
%The maximal chains representing $\ParS{\mathcal{S}}$ are $c-a-b$ and
%$c-d$. Every member of ${\mathcal{S}}$ has a match in
%$\Proj{\mathcal{S}}$, and vice versa.
The set $\mathcal{S}$ is order-closed,
but, as will be explained below, it is not realisable by a single
DTS. Therefore, $\mathcal{S}$ must be partitioned, somehow.
The pairs of events not related by $\ParS{\mathcal{S}}$ are $(a,d)$ and
$(b,d)$, based on which two different OBDs,
 $(\{\SeqS_2,\SeqS_3\}$, $\{\SeqS_1\})$ and $(\{\SeqS_1,\SeqS_2\}$,
$\{\SeqS_3\})$, could be formed.
The divisions have not only equal o-rankings, but also equal
K-rankings and l-rankings.
So a division must be chosen randomly. Choosing the second division
results in the trivial expression
$\SeqS_1\BAR\SeqS_2\BAR\SeqS_3$. But, in
fact, a satisfactory result can be obtained from the first division.
\end{example}
\noindent
This example shows that in the presence of non-default constraints, when
${\mathcal{S}}$ is order-closed, but not realisable by a single DTS,
reasoning based on orderings alone
does not provide us with enough information to properly partition
${\mathcal{S}}$.
Next, we investigate how the constraints can be used to effectively
partition $\mathcal{S}$.
%\emph{We assume $\mathcal{S}$ is order-closed:
%it is not possible to divide it based on orderings.}
%---------------------
\begin{definition}
\label{def-locally-stable}
Let $\mathcal{S}$ be an order-closed set of scenarios.
%$\DR{\mathcal{S}}$=(\ParS{\mathcal{S}},\DF{\mathcal{S}})$
%be its distance relation. Let
Let $\SeqSC\in\mathcal{S}$ and
%$\DF{\SeqSC}$ be its distance function. Let
$p\in\Proj{\mathcal{S}}$, such that $\SeqSC$ matches $p$.
%$match(\SeqSC,p)$.
%generated from $\DR{\mathcal{S}}$, such that $\ESeq{\SeqSC}=\ESeq{p}$.
We say $\SeqSC$ is \emph{compatible with} $\DR{\mathcal{S}}$
if $\DF{\SeqSC}=\DF{p}$.
\end{definition}
%-------------------
\noindent
%Recall that projection $p$ is obtained from $\DR{\mathcal{S}}$
%(see \Sec{sec-dts}).
%So the notion of compatibility of a
%scenario in $\mathcal{S}$ depends on $\DR{\mathcal{S}}$.
Notice that if all the members of $\mathcal{S}$ have
only default constraints, then they are all compatible with
$\DR{\mathcal{S}}$.

In Example \ref{ex-with-time-1}
$\PR{\SeqS_1}$ is represented by $c-d-a-b$, while $\PR{\SeqS_2}$ and
$\PR{\SeqS_3}$ are represented by $c-a-d-b$ and $c-a-b-d$, respectively.
The distance functions of members of $\mathcal{S}$ are shown below:

$\DF{\SeqS_1}=
\{
(a, b, 0, 3),
 (c, a, 0, 3),
(c, b, 0, 3),
 (c, d, 0, 3),
(d, a, 0, 3),
(d, b, 0, 3)
\}$,

$\DF{\SeqS_2}=
\{
(a, b, 0, 3),
(a, d, 0, 1),
(c, a, 0, 1),
(c, b, 0, 3),
(c, d, 0, 1),
(d, b, 0, 3)\}$,

$\DF{\SeqS_3}=
\{
(a, b, 0, 1),
(a, d, 0, 1),
(b, d, 0, 1),
(c, a, 0, 1),
(c, b, 0, 1),
(c, d, 0, 1)
\}$.

The stable distance relation
of $\mathcal{S}$ (see Def. \ref{def-DR-S})
is $\DR{\mathcal{S}}=(\ParS{\mathcal{S}},\DF{\mathcal{S}})$, where
$\ParS{\mathcal{S}}=\PR{\SeqS_1}\cap\PR{\SeqS_2}\cap\PR{\SeqS_3}$ is
represented by $c-a-b$ and $c-d$, and
 $\DF{\mathcal{S}}$ is represented by
 $\{
(a, b, 0, 3),
(c, a, 0, 3),
(c, b, 0, 3),
(c, d, 0, 3)\}$.
Observe that, for example, $\DF{\mathcal{S}}(a,b)=[0,3]$ is the
smallest interval that includes $\DF{\SeqS_1}(a,b)$,
$\DF{\SeqS_2}(a,b)$ and $\DF{\SeqS_3}(a,b)$.
The set of projections of $\DR{\mathcal{S}}$ (see \Sec{sec-dts}) is
$\Proj{\mathcal{S}}=\{p_1,p_2,p_3\}$, where
$p_1$, $p_2$ and $p_3$ match $\SeqS_1$, $\SeqS_2$ and
$\SeqS_3$, respectively.
The distance functions of the projections in $\Proj{\mathcal{S}}$ are
shown below:

$\DF{p_1}=
\{
(a, b, 0, 3),
 (c, a, 0, 3),
(c, b, 0, 3),
 (c, d, 0, 3),
(d, a, 0, 3),
(d, b, 0, 3)
\}$,

$\DF{p_2}=
\{
(a, b, 0, 3),
(a, d, 0, 3),
(c, a, 0, 3),
(c, b, 0, 3),
(c, d, 0, 3),
(d, b, 0, 3)\}$,

$\DF{p_3}=
\{
(a, b, 0, 3),
(a, d, 0, 3),
(b, d, 0, 3),
(c, a, 0, 3),
(c, b, 0, 3),
(c, d, 0, 3)
\}$.

\noindent
Notice that $\DF{\SeqS_1}=\DF{p_1}$, so $\SeqS_1$ is compatible with
$\DR{\mathcal{S}}$, but $\SeqS_2$ and $\SeqS_3$ are not.
This observation can be immediately used to divide $\mathcal{S}$
into
%two partitions,
%, one that includes all members of $\mathcal{S}$ that are
%compatible with it, and one that includes the rest:
$\mathcal{S}_1=\{\SeqS_1\}$ and
$\mathcal{S}_2=\{\SeqS_2,\SeqS_3\}$.
The set $\mathcal{S}_2$ is order-closed,
%so we construct
its distance relation is
$\DR{\mathcal{S}_2}=(\ParS{\mathcal{S}_2},\DF{\mathcal{S}_2})$, where
$\ParS{\mathcal{S}_2}$ is represented by the maximal chains $c-a-b$
and $c-a-d$, and $\DF{\mathcal{S}_2}$ is represented by
 $\{
(a, b, 0, 3),
(a, d, 0, 1),
(c, a, 0, 1),
(c, b, 0, 3),
(c, d, 0, 1)\}$.
The set of projections of $\DR{\mathcal{S}_2}$,
$\Proj{\mathcal{S}_2}$, includes two projections, $p'_2$ and $p'_3$,
that
are identical to $\SeqS_2$ and $\SeqS_3$, respectively, so
$\DF{\SeqS_2}=\DF{p'_2}$ and $\DF{\SeqS_3}=\DF{p'_3}$:
%they  are identical to their respective
%matching projections,
$\SeqS_2$ and $\SeqS_3$ are compatible with $\DR{\mathcal{S}_2}$.
The set $\mathcal{S}_2$ is realised by DTS $\xi_1||\xi_2$, where
$\xi_1=(cab,\{\TT{c}{a}\leq 1,\TT{c}{b}\leq
3\})$ and $\xi_2=(cad,\{\TT{c}{d}\leq 1\})$.
So ${\mathcal{S}}$ is realised by $\SeqS_1\BAR(\xi_1||\xi_2)$
(compare this with the trivial solution shown in Example \ref{ex-with-time-1}).
%the expression $\D(\SeqS_1)\BAR E$.
%, where $E_2=\{\SeqS_1\}$.
%--------------------------
%\vspace{-0.5em}
%\begin{obs}
%\label{obs-realisable}
%An order-closed set of scenarios, $\mathcal{S}$, is realisable by a DTS iff
%all its members are compatible with $\DR{\mathcal{S}}$.
%\end{obs}
%\noindent

In Example \ref{ex-with-time-1} we partitioned $\mathcal{S}$ into
scenarios that are compatible with $\DR{\mathcal{S}}$, and those that
are not. But sometimes none of the members of $\mathcal{S}$ will be
compatible with $\DR{\mathcal{S}}$, so we need another criterion for
partitioning.
%---------------------------
%\vspace{-0.5em}
\begin{definition}
\label{def-c-dis}
Let $\mathcal{S}$ be an order-closed set of scenarios and
$\SeqSC\in\mathcal{S}$.
%and $p\in\Proj{\mathcal{S}}$,
%such that $\SeqSC$ matches $p$.
%and $match(\SeqSC,p)$.
The \emph{constraint-based disagreement} between $\SeqSC$ and
$\DR{\mathcal{S}}$
%, denoted by $\CDis{\mathcal{S}}(\SeqSC,p)$,
is $\CDis{\mathcal{S}}(\SeqSC)$=
\Set{(e,e')}{e\PR{\SeqSC} e'\wedge
  \DF{\SeqSC}(e,e')\neq\DF{p}(e,e') \text{ where }
  p \text{ is the projection of } \DR{\mathcal{S}} \text{ that matches }
  \SeqSC}.
\end{definition}
%-----------------------
%\vspace{-1em}
\begin{definition}
\label{def-CBD}
Let $\mathcal{S}$ be an order-closed set of scenarios over
$\Sigma_{\mathcal{S}}$, such that none of its members are compatible
with $\DR{\mathcal{S}}$ and $|\mathcal{S}|>1$. Let $e,e'\in\Sigma_{\mathcal{S}}$.
The \emph{constraint-based division} (\emph{CBD} for short) \emph{of}
$\mathcal{S}$ \emph{with respect to} $(e,e')$, denoted by
$\CBD{\mathcal{S}}(e,e')$, is the unordered pair
$(\mathcal{S}_1,\mathcal{S}_2)$, such that
\setlist{nolistsep}
\begin{enumerate}[(1), noitemsep]
 % $(e,e')\in\CDis{\mathcal{S}}(\SeqS_i,p_i)$ for some $\SeqS_i\in\mathcal{S}$
  \item
    $\mathcal{S}_1=$\Set{\SeqSC\in\mathcal{S}}{(e,e')\in\CDis{\mathcal{S}}(\SeqSC)}
    %{e\PR{\SeqSC}e'\wedge
  %\DF{\SeqSC}(e,e')\neq\DF{p}(e,e'), \text{ where }
    %p\in\Proj{\mathcal{S}}\wedge match(\SeqSC,p)}
    and
    $\mathcal{S}_2=\mathcal{S}\setminus\mathcal{S}_1$, or
  \item
    $(e,e')\not\in\CDis{\mathcal{S}}(\SeqSC)$ for any
    $\SeqSC\in\mathcal{S}$ and there exists some $\SeqSC'\in\mathcal{S}$, where
    $e\PR{\SeqSC'}e'\wedge\DF{\SeqSC'}(e,e')\neq[0,\infty]$,
    %or $e'\PR{\SeqSC'}e\wedge\DF{\SeqSC'}(e',e)\neq[0,\infty]$,
    then
    $\mathcal{S}_1=$\Set{\SeqS\in\mathcal{S}}{e\ParS{\SeqS}e'} and
     $\mathcal{S}_2=\mathcal{S}\setminus\mathcal{S}_1$.
   % $\mathcal{S}_2=$\Set{\SeqS\in\mathcal{S}}{e'\ParS{\SeqS}e}.
    \end{enumerate}
  %  We say members of $\mathcal{S}_1$ and $\mathcal{S}_2$
  %  \emph{disagree} on $(e,e')$.
\end{definition}
%\vspace{-0.5em}
%-------------------
\noindent
Intuitively, if
$\CBD{\mathcal{S}}(e,e')=(\mathcal{S}_1,\mathcal{S}_2)$ is formed by case (1)
of Def. \ref{def-CBD}, then all the members of $\mathcal{S}_1$
disagree with their respective
matching projections on the interval for $(e,e')$, i.e.,
on $\DF{}(e,e')$.
In this case, if $\SeqS_2\in\mathcal{S}_2$, then either
$\DF{\SeqS_2}(e,e')=\DF{p_2}(e,e')$, where $p_2$ is the
projection that matches $\SeqS_2$, or $e\not\PR{\SeqS_2}e'$.
Notice that $\mathcal{S}_1\neq\emptyset$, but
$\mathcal{S}_2$ can be empty. If $\mathcal{S}_2=\emptyset$, then
$(\mathcal{S}_1,\mathcal{S}_2)$ is \emph{not} a CBD.
%If $\mathcal{S}_2\neq\emptyset$, we call
%$(\mathcal{S}_1,\mathcal{S}_2)$ a \emph{non-trivial} division of
%$\mathcal{S}$.

If
$\CBD{\mathcal{S}}(e,e')=(\mathcal{S}_1,\mathcal{S}_2)$ is formed by case (2),
then all the members of $\mathcal{S}_1$ agree with
their respective matching projections on the interval for $(e,e')$, however,
at least one of them has a non-default interval on $(e,e')$.
%-------------
%Feliks-almost-gap
%\vspace{-0.5em}
\begin{example}
\label{ex-with-time-2}
Let $\mathcal{S}=\{\SeqS_1,\SeqS_2\}$, where
 %over $\Sigma=\{a,b,c,d\}$, where
$\SeqS_1=(abc,\{1\leq\TT{a}{b}\leq 2\})$ and
$\SeqS_2=(acb,\{2\leq\TT{a}{b}\leq 3\})$.
The partial order in $\mathcal{S}$,
$\ParS{\mathcal{S}}=\PR{\SeqS_1}\cap\PR{\SeqS_2}$, is
represented by $a-b$ and $a-c$.
The distance functions of the members of $\mathcal{S}$ are
$\DF{\SeqS_1}=
\{
(a, b, 1, 2),
 (a, c, 1, \infty)
\}$ and
$\DF{\SeqS_2}=
\{
(a, b, 2, 3),
 (a, c, 0, 3),
  (c, b, 0, 3)
\}$.

Observe that the smallest interval that includes
$\DF{\mathcal{\SeqS}_1}(a,b)=[1,2]$ and
$\DF{\mathcal{\SeqS}_2}(a,b)=[2,3]$ is $[1,3]$,
so $\DF{\mathcal{S}}(a,b)=[1,3]$.
Similarly,
$\DF{\mathcal{S}}(a,c)=[0,\infty]$ is the smallest interval that includes
$\DF{\mathcal{\SeqS}_1}(a,c)=[1,\infty]$ and
$\DF{\mathcal{\SeqS}_2}(a,c)=[0,3]$. Therefore,
$\DF{\mathcal{S}}$ is represented by $\{(a, b, 1, 3)\}$ (recall that
we do not show the default intervals).
%The distance relation of $\mathcal{S}$,
%$\DR{\mathcal{S}}=(\ParS{\mathcal{S}},\DF{\mathcal{S}})$, where
%$\ParS{\mathcal{S}}=\{(a,b),(a,c)\}$ and
%maximal chains, $a-b$ and $a-c$,

The projections of $\DR{\mathcal{S}}$ yield
 $\DF{p_1}=
 \{
 (a, b, 1, 3),
  (a, c, 1, \infty)
\}$
 and
$\DF{p_2}=
\{
 (a,b,1,3),
 (a,c,0,3),
 (c, b, 0, 3)
\}$, where $\SeqS_1$ matches $p_1$ and $\SeqS_2$ matches $p_2$.
%$match(\SeqS_1,p_1)$ and $match(\SeqS_2,p_2)$.
None of the members of $\mathcal{S}$ are compatible with
$\DR{\mathcal{S}}$:
$\CDis{\mathcal{S}}(\SeqS_1)=\CDis{\mathcal{S}}(\SeqS_2)=\{(a,
b)\}$.
By case (1) of Def. \ref{def-CBD},
$\CBD{\mathcal{S}}(a,b)=(\{\SeqS_1,\SeqS_2\},\emptyset)$.
%Observe that $\DF{\SeqS_1}(a,b)\neq\DF{p_1}(a,b)$ and
%$\DF{\SeqS_2}(a,b)\neq\DF{p_2}(a,b)$, so
However,
$\DF{\SeqS_1}(b,c)=[0,\infty]$ and
$\DF{\SeqS_2}(c,b)=[0,3]\neq[0,\infty]$, so, by case (2) of
Def. \ref{def-CBD},
$\CBD{\mathcal{S}}(b,c)=(\{\SeqS_1\},\{\SeqS_2\})$.
%Observe that $\CBD{\mathcal{S}}(c,b)=(\{\SeqS_2\},\{\SeqS_1\})$,
\end{example}
%------------------
% new-crit-2
%\vspace{-0.5em}
%\begin{example}
%\label{ex-with-time-4}
%Let $\mathcal{S}=\{\SeqS_1,\SeqS_2,\SeqS_3\}$, where
%$\SeqS_1=(adcb,\{\TT{d}{c} \leq 2,\TT{a}{b} \geq 3\})$,
%$\SeqS_2=(dcab,\{\TT{d}{c} \geq 1,\TT{d}{a} \geq 2\})$ and
%$\SeqS_3=(dacb,\{\TT{d}{a} \geq 2\})$.
%%The distance functions of the members of $\mathcal{S}$ are
%
%\noindent
%$\DF{\SeqS_1}=
% \{
%  (a, b, 3, \infty),
%(d, c, 0, 2)
%\}$,
%$\DF{\SeqS_2}=
%\{
% (d,a,  2, \infty),
%  (d, b, 2, \infty),
% (d, c, 1, \infty)
% \}$ and
%$\DF{\SeqS_3}=
% \{
% (d, a, 2, \infty),
%  (d, b, 2, \infty),
%(d,c,  2, \infty)
% \}$.
%%$\DR{\mathcal{S}}=(\ParS{\mathcal{S}},\DF{\mathcal{S}})$, where
%$\ParS{\mathcal{S}}$ is represented by $d-c-b$ and $a-b$, and
%$\DF{\mathcal{S}}$ has only default intervals.
% The projections of $\DR{\mathcal{S}}$ are
% %that match the members of
%%$\mathcal{S}$ are, respectively,
%$p_1=(adcb,\emptyset)$,
%$p_2=(dcab,\emptyset)$ and
%$p_3=(dacb,\emptyset)$. Observe that
%$\mathcal{S}$ is order-closed, but none of its members are compatible
%with $\DR{\mathcal{S}}$.
%%So, by Obs. \ref{obs-realisable}, $\mathcal{S}$
%%is not realisable by a single DTS.
%We have
%
%\noindent
%$\CDis{\mathcal{S}}(\SeqS_1)=\{(a, b),(d, c)\}$,
%$\CDis{\mathcal{S}}(\SeqS_2)=\CDis{\mathcal{S}}(\SeqS_3)=\{
%(d, c),(d,a),(d, b)\}$.
%%\noindent
%%$\CDis{\mathcal{S}}(\SeqS_3,p_3)=\{ (d, a),(d,c),(d, b)\}$.
%
%\noindent
%%By case (1) of \Def{def-CBD},
%$\CBD{\mathcal{S}}(a,b)=(\{\SeqS_1\},\{\SeqS_2,\SeqS_3\})$,
%$\CBD{\mathcal{S}}(d,a)=\CBD{\mathcal{S}}(d,b)=(\{\SeqS_2,\SeqS_3\},\{\SeqS_1\})$
%and
%
%\noindent
%$\CBD{\mathcal{S}}(d,c)=(\{\SeqS_1,\SeqS_2,\SeqS_3\},\emptyset)$.
%%\noindent
%%$\CBD{\mathcal{S}}(d,b)=(\{\SeqS_2,\SeqS_3\},\{\SeqS_1\})$.
%While $\CBD{\mathcal{S}}(d,c)$ is not so useful, the other CBDs
%effectively divide $\mathcal{S}$ into
%%partitions
%$\mathcal{S}_1=\{\SeqS_2,\SeqS_3\}$ and $\mathcal{S}_2=\{\SeqS_1\}$,
%which are both
%%Set $\mathcal{S}_1$ is
%order-closed.
%$\DR{\mathcal{S}_1}=(\ParS{\mathcal{S}_1},\DF{\mathcal{S}_1})$, where
%$\ParS{\mathcal{S}_1}$ is represented by $d-a-b$ and $d-c-b$, while
%%$\ParS{\mathcal{S}_1}=\{(a,b),(d,c),(c,b),(d,b),(d,a)\}$ and
% $\DF{\mathcal{S}_1}$ is represented by
% $\{
%(d, a, 2, \infty),
%(d, b, 2, \infty),
%(d, c, 1, \infty)\}$.
% The projections of $\DR{\mathcal{S}_1}$ are identical to
% the respective members of $\mathcal{S}_1$, so
% $\SeqS_2$ and $\SeqS_3$ are compatible with $\DR{\mathcal{S}_1}$.
%%By Obs. \ref{obs-realisable}, $\mathcal{S}_1$ is
%%realisable:
%The DTS $\xi_1||\xi_2$, where $\xi_1=(dab,\{\TT{d}{a}\geq 2\})$ and
%$\xi_2=(dcb,\{\TT{d}{c}\geq 1,\TT{d}{b}\geq 2\})$ realises
%$\mathcal{S}_1$. So ${\mathcal{S}}$ is realised by
%%the expression
%$(\xi_1||\xi_2)\BAR\SeqS_1$.
%\end{example}
%-------------------------
%THE THEOREM USED TO BE HERE !!!!
%\vspace{-0.1em}
\begin{theorem}
\label{thm-not-locally-stable}
Let $\mathcal{S}$ be an order-closed set of scenarios, such that
none of its members is compatible with $\DR{\mathcal{S}}$
%$|\ESeq{\SeqSC}|>3$, for each $\SeqSC\in\mathcal{S}$,
and $|\mathcal{S}|>1$.
Then there is at least one pair of events $e$ and $e'$ such that
$\CBD{\mathcal{S}}(e,e')=(\mathcal{S}_1,\mathcal{S}_2)$, where
$\mathcal{S}_1\neq\emptyset$ and $\mathcal{S}_2\neq\emptyset$.
\end{theorem}
%\vspace{-0.1em}
%\begin{proof}
%The proof is presented in Appendix A.
% \end{proof}
%---------------------------
\noindent
In general, $\mathcal{S}$
might have more than one
CBD, so we need a criterion for comparing them.
%in which case the one with the lowest K-ranking must be
%selected. But to break a tie
%we must find a new criterion based on the constraints of the
%partitions of CBDs.
%---------------------------
%--------------------
%\vspace{-0.1em}
\begin{definition}
  \label{def-c-distance}
  Let $\mathcal{S}$ be an order-closed set of scenarios and
  $\SeqS_i,\SeqS_j\in\mathcal{S}$ ($i\neq j$).
 % be two scenarios over the same set of events in $\mathcal{S}$.
  The \emph{c-distance} of $\SeqS_i$ and $\SeqS_j$,
 % with respect to $\mathcal{S}$,
%  denoted by $\CD{\SeqS_i}{\SeqS_j}{\mathcal{S}}$, is
   denoted by $\CD{\SeqS_i}{\SeqS_j}$, is
$|\Set{\CBD{\mathcal{S}}(e,e')=(\mathcal{S}_1,\mathcal{S}_2)}
    {\mathcal{S}_1\neq\emptyset\wedge\mathcal{S}_2\neq\emptyset\wedge
  %  e\PR{\SeqS_i} e'\wedge e\PR{\SeqS_j} e'\wedge
      ((\SeqS_i\in\mathcal{S}_1\wedge\SeqS_j\in\mathcal{S}_2)\vee
 (\SeqS_i\in\mathcal{S}_2\wedge\SeqS_j\in\mathcal{S}_1))}|$.
%$\K{\SeqS_1}{\SeqS_2}$, is the number of pairwise
%disagreements between $\ESeq{\SeqS_1}$ and $\ESeq{\SeqS_2}$.
\end{definition}
\noindent
Intuitively, the c-distance between $s_i$ and $s_j$ is the
number of times that $s_i$ and $s_j$ would wind up in different
partitions of CBDs.
The smaller the c-distance, the more similar are the constraints
of the scenarios.
%constraint-based divisions.
%\vspace{-0.1em}
\begin{definition}
\label{def-c-set-ranking}
Let $\mathcal{S}$ be an order-closed set of scenarios.
The \emph{c-ranking} of ${\mathcal{S}}$, denoted by $\CK{\mathcal{S}}$,
is the average of all pairwise c-distances between its members.
%the members of $\mathcal{S}$.
\end{definition}
%------------------
%\vspace{-0.1em}
\begin{definition}
\label{def-c-ranking}
Let $\mathcal{S}$ be an order-closed set of scenarios and
$D=(\mathcal{S}_1,\mathcal{S}_2)$ be a CBD of
$\mathcal{S}$. The \emph{c-ranking} of $D$, denoted by
$\CK{D}$, is given by
$\CK{\mathcal{S}_1}+\CK{\mathcal{S}_2}$.
%$\RK{\mathcal{S}_1}*\frac{|\mathcal{S}_1|}{|\mathcal{S}|}+
%\RK{\mathcal{S}_2}*\frac{|\mathcal{S}_2|}{|\mathcal{S}|}$.
\end{definition}
%\noindent
%A low c-ranking of a CBD is an indication of similarity of the
%constraints between members of the partitions in the CBD.
% unknown2
%\vspace{-0.1em}
\begin{example}
\label{ex-with-time-3}
Consider the set $\mathcal{S}=\{\SeqS_1,\SeqS_2,\SeqS_3\}$, where
 %over $\Sigma=\{a,b,c,d\}$, where
$\SeqS_1=(acdb,\{\TT{d}{b} \geq 2\})$,
$\SeqS_2=(adbc,\{\TT{d}{c} \geq 2\})$ and
$\SeqS_3=(adcb,\{\TT{d}{b} \geq 2\})$.
The set $\mathcal{S}$ is order-closed, but not realisable by a DTS.
The distance functions of its members are
$\DF{\SeqS_1}=
 \{
% (a, d, 0, \infty),
% (a,c,  0, \infty),
  (a, b, 2, \infty),
(c, b, 2, \infty),
(d, b, 2, \infty)
%(a, b, 0, \infty)
\}$,
$\DF{\SeqS_2}=
 \{
 (a, c, 2, \infty),
%   (d, b, 2, \infty),
 (d,c,  2, \infty)
%(c, a , 0, \infty),
%(c, b, 0, \infty),
 %(a, b, 0, \infty)
 \}$ and
$\DF{\SeqS_3}=
 \{
  (a, b, 2, \infty),
% (d,c,  2, \infty),
  (d, b, 2, \infty)
%(a, c , 0, \infty),
%(a , b, 0, \infty),
 %(c, b, 0, \infty)
 \}$.

 %$\DR{\mathcal{S}}=(\ParS{\mathcal{S}},\DF{\mathcal{S}})$, where
The partial order in $\mathcal{S}$, $\ParS{\mathcal{S}}$, is represented by
$a-d-b$ and $a-c$, while $\DF{\mathcal{S}}$ has only default intervals.
% The members of $\Proj{\mathcal{S}}$ are
The projections of $\DR{\mathcal{S}}$ are
$p_1=(acdb,\emptyset)$,
$p_2=(adbc,\emptyset)$ and
$p_3=(adcb,\emptyset)$.
So none of the members of $\mathcal{S}$ are compatible with
$\DR{\mathcal{S}}$.

The constraint-based disagreement between members of $\mathcal{S}$ and
$\DR{\mathcal{S}}$ are

\noindent
$\CDis{\mathcal{S}}(\SeqS_1)=\{(a,b),(c,b),(d,b)\}$,
$\CDis{\mathcal{S}}(\SeqS_2)=\{(a,c),(d,c)\}$ and
$\CDis{\mathcal{S}}(\SeqS_3)=\{(a,b),$\\$(d,b)\}$,
while
$\CBD{\mathcal{S}}(a,b)=\CBD{\mathcal{S}}(d,b)=(\{\SeqS_1,\SeqS_3\},\{\SeqS_2\})$,
$\CBD{\mathcal{S}}(a,c)=\CBD{\mathcal{S}}(d,c)=(\{\SeqS_2\},\{\SeqS_1,\SeqS_3\})$ and
$\CBD{\mathcal{S}}(c,b)=(\{\SeqS_1\},\{\SeqS_2,\SeqS_3\})$.
%$\CBD{\mathcal{S}}(d,b)=(\{\SeqS_1,\SeqS_3\},\{\SeqS_2\})$,
%$\CBD{\mathcal{S}}(d,c)=(\{\SeqS_2\},\{\SeqS_1,\SeqS_3\})$.

\noindent
This analysis leads to divisions $D_1=(\{\SeqS_2\},\{\SeqS_1,\SeqS_3\})$
and $D_2=(\{\SeqS_1\},\{\SeqS_2,\SeqS_3\})$, which have
equal o-rankings and K-rankings,
so we compare their c-rankings:
\noindent
$\CK{\{\SeqS_2\}}= 0$,
$\CK{\{\SeqS_1,\SeqS_3\}}=\CD{\SeqS_1}{\SeqS_3}=|\{\CBD{\mathcal{S}}(c,b)\}|=1$,
$\CK{\{\SeqS_1\}}= 0$
and
$\CK{\{\SeqS_2,\SeqS_3\}}=\CD{\SeqS_2}{\SeqS_3}=
|\{\CBD{\mathcal{S}}(a,b),\CBD{\mathcal{S}}(a,c),\CBD{\mathcal{S}}(d,b),\CBD{\mathcal{S}}(d,c)\}|=4$.
So $\CK{D_1}=0+1=1$ and $\CK{D_2}=0+4=4$.

In $D_1$, $\{\SeqS_1,\SeqS_3\}$ is realisable by the DTS
$\xi_1||\xi_2$, where $\xi_1=(acb,\{\TT{a}{b}\geq 2\})$ and
$\xi_2=(adb,\{\TT{d}{b}\geq 2\})$.
So ${\mathcal{S}}$ is realised by $(\xi_1||\xi_2)\BAR\SeqS_2$.
If we choose $D_2$ (which has the higher c-ranking), even though
$\{\SeqS_2,\SeqS_3\}$ is order-closed, its members would not be
compatible with the distance relation of the set. So we would have to
divide the set once more and end up with the trivial expression
$\SeqS_1\BAR\SeqS_2\BAR\SeqS_3$.
\end{example}
%-----------------
%\begin{obs}
%  Let $\mathcal{S}$ be an order-closed set of scenarios.
%  All the members of $\mathcal{S}$ are compatible with
%  $\DR{\mathcal{S}}$ iff $\CK{\mathcal{S}}=0$.
%  \end{obs}
\noindent
As shown above, if none of the members of an order-closed set,
%of scenarios,
$\mathcal{S}$, are compatible with $\DR{\mathcal{S}}$, then we compute all the
CBDs of $\mathcal{S}$ and
%By Theorem \ref{obs-not-locally-stable}, there must be
%at least one non-trivial CBD.
choose the one with the lowest o-ranking.
We break a tie by choosing the CBD with the lowest
K-ranking.
If K-rankings cannot break a tie, we
choose a division with the lowest c-ranking, and if necessary we use
l-rankings to break a tie.
%whose partitions are order-closed.
%A division with two order-closed partitions has priority over one
%that has only one order-closed partition, and that has priority over
%one that has no order-closed partitions.
%If order-closedness of partitions cannot break a tie, we randomly
%choose a division, while avoiding partitions with gaps if possible.
%\vspace{-0.4em}
\begin{theorem}
\label{thm-realisable}
%Let $\mathcal{S}$ be
A finite set, $\mathcal{S}$, of sequential scenarios is
realisable by a single DTS iff it is order-closed and every member of
it is compatible with $\DR{\mathcal{S}}$.
\end{theorem}
%\vspace{-0.5em}
%\begin{proof}
%The proof is presented in Appendix B.
%\end{proof}
\noindent
We now present an algorithm
%%Next, we develop an algorithm
that solves the problem of realisability
of a set of scenarios as a scenario expression.
%%%%%%%%%%%%%%%%%%%%%%%


%%%%%%%%%%% GAP %%%%%%



%The next example shows that given a set of sequential scenarios there
%might be more than one expression that realises the set.
